% =====================================================================
% AI Academy -- National AI Center
% Software Engineering Final Project -- Team Report
%
% Team: Byte Bite  |  Topic: 2 -- AI Food Analyzer
% =====================================================================

\documentclass[11pt,a4paper]{article}

% ===== EARLY MACROS ==================================================
\usepackage[dvipsnames]{xcolor}
\definecolor{accent2}{HTML}{8B0000}
\definecolor{ph}{HTML}{6B7280}
\newcommand{\TODO}[1]{\textcolor{accent2}{\textbf{[TODO: #1]}}}
\newcommand{\phh}[1]{\textcolor{ph}{\textit{#1}}}

% ===== CONFIGURATION =================================================
\newcommand{\teamname}{Byte Bite}
\newcommand{\topicchoice}{Topic 2 --- AI Food Analyzer}
\newcommand{\member}[1]{\textbf{#1}}
\newcommand{\reportdate}{May 22, 2026}
\newcommand{\repourl}{https://github.com/AMammedova/food-analyzer-byte-bite}

% ===== COURSE META ===================================================
\newcommand{\coursecode}{AI-ENG-110}
\newcommand{\coursename}{Software Engineering}
\newcommand{\institution}{AI Academy, National AI Center}
\newcommand{\semester}{Spring 2026}

% ===== PACKAGES ======================================================
\usepackage[margin=2.2cm,headheight=15pt]{geometry}
\usepackage[T1]{fontenc}
\usepackage[utf8]{inputenc}
\usepackage[final,protrusion=true,expansion=false]{microtype}
\usepackage{amsmath,amssymb}
\usepackage{graphicx}
\usepackage{booktabs}
\usepackage{array}
\usepackage{tabularx}
\usepackage{longtable}
\usepackage{ragged2e}
\usepackage{enumitem}
\usepackage[parfill]{parskip}
\usepackage{listings}
\usepackage[most]{tcolorbox}
\usepackage{titlesec}
\usepackage{fancyhdr}
\usepackage{lastpage}
\usepackage[hidelinks]{hyperref}

\emergencystretch=3em
\hbadness=10000
\hfuzz=2pt

\hypersetup{
  colorlinks=true,
  linkcolor=NavyBlue,
  urlcolor=NavyBlue,
  citecolor=NavyBlue,
  pdftitle={\coursecode\ Final Project Report},
  pdfauthor={\teamname}
}

% ===== STYLING =======================================================
\definecolor{accent}{HTML}{0B3D91}
\definecolor{soft}{HTML}{F2F4F8}
\definecolor{archbg}{HTML}{F8F9FA}

\titleformat{\section}
  {\Large\bfseries\color{accent}}{\thesection.}{0.6em}{}
\titleformat{\subsection}
  {\large\bfseries\color{accent}}{\thesubsection}{0.5em}{}
\titleformat{\subsubsection}
  {\normalsize\bfseries\color{accent!80!black}}{\thesubsubsection}{0.4em}{}

\let\ph\phh

% ===== CODE STYLE ====================================================
\definecolor{codegreen}{rgb}{0,0.55,0}
\definecolor{codegray}{rgb}{0.4,0.4,0.4}
\definecolor{codepurple}{rgb}{0.55,0,0.78}
\definecolor{backcolour}{rgb}{0.97,0.97,0.97}

\lstdefinestyle{python}{
  language=Python,
  backgroundcolor=\color{backcolour},
  commentstyle=\color{codegreen},
  keywordstyle=\color{accent},
  numberstyle=\tiny\color{codegray},
  stringstyle=\color{codepurple},
  basicstyle=\ttfamily\footnotesize,
  breakatwhitespace=false,
  breaklines=true,
  keepspaces=true,
  numbers=left,
  numbersep=6pt,
  tabsize=2,
  frame=single,
  framerule=0.4pt,
  rulecolor=\color{black!30},
}
\lstdefinestyle{bash}{
  language=bash,
  backgroundcolor=\color{backcolour},
  basicstyle=\ttfamily\small,
  breaklines=true,
  frame=single,
  framerule=0.4pt,
  rulecolor=\color{black!30},
  numbers=none,
}
\lstdefinestyle{sql}{
  language=SQL,
  backgroundcolor=\color{backcolour},
  keywordstyle=\color{accent},
  basicstyle=\ttfamily\small,
  breaklines=true,
  frame=single,
  framerule=0.4pt,
  rulecolor=\color{black!30},
  numbers=none,
}

% ===== PAGE STYLE ====================================================
\pagestyle{fancy}
\fancyhf{}
\fancyhead[L]{\small\textit{\coursecode\ -- Final Report -- \teamname}}
\fancyhead[R]{\small\textit{\semester}}
\fancyfoot[L]{\small\textit{\institution}}
\fancyfoot[R]{\small Page \thepage\ of \pageref{LastPage}}
\renewcommand{\headrulewidth}{0.4pt}
\renewcommand{\footrulewidth}{0.4pt}

% =====================================================================
\begin{document}

% ===== TITLE BLOCK ===================================================
\begin{center}
{\small \textsc{\institution} \\[0.2em]
 \textsc{\coursecode\ -- \coursename\ -- \semester}}\\[1.2em]
{\Large\bfseries\color{accent} Final Project Report}\\[0.3em]
{\LARGE\bfseries \topicchoice}\\[0.8em]
\rule{0.85\textwidth}{0.6pt}
\end{center}

\vspace{0.5em}
\begin{tabularx}{\textwidth}{@{}lX lX@{}}
\textbf{Team:}  & \teamname          & \textbf{Date:} & \reportdate \\
\textbf{Tag:}   & \texttt{v1.0-final}& \textbf{Repo:} & \small\url{\repourl} \\
\end{tabularx}

\smallskip
\noindent\textbf{Members:}\enspace
\member{Aysel Mamedova},\enspace
\member{Gülnur Mammadova},\enspace
\member{Şamistan Hüseynov},\enspace
\member{Rahima Karimova}

\vspace{0.4em}\hrule\vspace{0.6em}

% ===== EXECUTIVE SUMMARY =============================================
\section*{Executive Summary}
\addcontentsline{toc}{section}{Executive Summary}

We built \textbf{Food Analyzer (Byte Bite)}, a production-grade web service that accepts a meal photograph, identifies ingredients and their estimated portion sizes via a Google Gemini vision-language model, retrieves per-ingredient nutrition data from the USDA FoodData Central API, and returns a full macronutrient breakdown (kcal, protein, carbs, fat).
The headline concurrency result is a \textbf{4.8$\times$} wall-clock speedup when fetching nutrition data for five ingredients concurrently rather than sequentially (2.50\,s $\to$ 0.53\,s), achieved with \texttt{asyncio.gather} bounded by a \texttt{Semaphore(10)}.
The most important robustness story is our multi-layer API-failure defence: every call to the Gemini VLM and USDA endpoints is wrapped with \texttt{tenacity} exponential backoff (4 attempts, 2--30\,s), a per-call timeout, an in-memory TTL cache, and structured logging.
The single most surprising lesson was that \texttt{pydantic-settings} intentionally does \emph{not} write \texttt{.env} values into \texttt{os.environ}, which broke provider adapters that relied on \texttt{os.getenv} directly; we fixed this with a targeted \texttt{load\_dotenv} call in \texttt{api.py}.

\tableofcontents
\vspace{0.6em}\hrule\vspace{0.4em}

% =====================================================================
\section{Project Overview}

\subsection{Problem framing \& scope}

The AI Food Analyzer solves a practical everyday problem: a person photographs a meal and wants to know its nutritional content without manually looking up each ingredient.
The system accepts any JPEG or PNG image (up to 5\,MB), passes it through a vision-language model (VLM) to extract a structured list of ingredients with estimated gram weights and confidence scores, then concurrently queries the USDA FoodData Central database for per-ingredient nutrition facts, and finally aggregates the results into a total calorie and macronutrient report.

\textbf{Entry points.} Users can interact through three channels:
\begin{enumerate}[topsep=2pt,itemsep=1pt]
  \item A browser-based drag-and-drop UI at \texttt{GET /}
  \item A REST API at \texttt{POST /analyze} (history at \texttt{GET /history})
  \item A CLI via \texttt{python -m foodanalyzer analyze <image\_path>}
\end{enumerate}
All three share the same underlying analysis pipeline in \texttt{core/analyzer.py}.

\textbf{Provider stack.}
We used Google Gemini (\texttt{gemini-2.5-flash}) for VLM ingredient identification and the USDA FoodData Central API for nutrition lookup.
The USDA key is free and requires only a sign-up at \texttt{api.data.gov}.

\textbf{Scope decisions.}
We implemented everything in the Topic~2 specification: HTTP API, CLI, PostgreSQL persistence, concurrent nutrition pipeline, retries, caching, validation, logging, and Docker.
Multi-provider failover was listed as an optional bonus; we focused instead on achieving high test coverage (79\,\%) and a polished web UI for demonstration.

\subsection{Provider choice and the AI module contract}

We chose \textbf{Google Gemini} (\texttt{LLM\_PROVIDER=gemini}, model \texttt{gemini-2.5-flash}) because the team had an active Google API key and the model performed well on food image recognition in early tests.
For nutrition data we used the provided \textbf{USDA FoodData Central} adapter (\texttt{ai.nutrition.USDAProvider}).

The AI-layer functions our code calls are:
\begin{itemize}[topsep=2pt,itemsep=1pt]
  \item \texttt{ai.identify\_ingredients(image\_path, vlm)}
  \item \texttt{ai.nutrition.USDAProvider.lookup(name)}
  \item \texttt{ai.calculator.compute\_totals(ingredients, facts)}
\end{itemize}

We did \textbf{not} modify the public interface of the provided \texttt{ai/} package; all changes are in wrappers and service layers under \texttt{src/foodanalyzer/}.
When the model returns semantically wrong output (negative \texttt{estimated\_grams}, confidence outside $[0,1]$, or empty ingredient list), our \texttt{FoodAnalyzer.analyze()} catches the case via Pydantic validation and either raises a descriptive \texttt{ValueError} or returns an empty result rather than propagating an exception.

% =====================================================================
\section{Architecture}\label{sec:arch}

\subsection{Module map}

\begin{tcolorbox}[colback=archbg, colframe=black!25, boxrule=0.4pt,
  arc=2pt, left=6pt, right=6pt, top=4pt, bottom=4pt,
  fontupper=\ttfamily\small]
Browser / curl\\
\phantom{xx}|\\
\phantom{xx}v\\
FastAPI  src/foodanalyzer/api.py\\
\phantom{xx}|-- POST /analyze\\
\phantom{xx}|     |-- validate\_image\_bytes()           [validation.py]\\
\phantom{xx}|     |-- AIService.identify\_ingredients()  [services/ai\_service.py]\\
\phantom{xx}|     |       |-- tenacity retry + 60 s timeout\\
\phantom{xx}|     |       `-- ai.identify\_ingredients() [PROVIDED ai/vlm.py]\\
\phantom{xx}|     |-- NutritionPipeline.fetch\_all()    [concurrency/pipeline.py]\\
\phantom{xx}|     |       |-- asyncio.gather + Semaphore(MAX\_PARALLEL)\\
\phantom{xx}|     |       |-- NutritionCache (TTL 24 h) [services/nutrition\_cache.py]\\
\phantom{xx}|     |       `-- ai.USDAProvider.lookup()  [PROVIDED ai/nutrition.py]\\
\phantom{xx}|     |-- ai.calculator.compute\_totals()   [PROVIDED ai/calculator.py]\\
\phantom{xx}|     `-- repository.save()                [storage/repository.py]\\
\phantom{xx}|             `-- asyncpg pool             [storage/db.py] --> PostgreSQL\\
\phantom{xx}|\\
\phantom{xx}`-- GET /history  --> repository.list\_recent() --> PostgreSQL\\
\\
CLI  src/foodanalyzer/cli.py  (shares core/analyzer.py)
\end{tcolorbox}

\subsection{Module summary}

\begin{center}\small
\begin{tabularx}{\textwidth}{@{} l X @{}}
\toprule
\textbf{Module} & \textbf{Responsibility} \\
\midrule
\texttt{config.py}                  & \texttt{pydantic-settings} \texttt{Settings} class; reads \texttt{.env}; singleton via \texttt{get\_settings()} \\
\texttt{models.py}                  & SE-layer Pydantic models: \texttt{IngredientOut}, \texttt{TotalsOut}, \texttt{AnalysisResult}, \texttt{AnalysisRecord} \\
\texttt{validation.py}              & Magic-byte format check (JPEG/PNG), configurable size limit \\
\texttt{logging\_config.py}         & \texttt{StructuredFormatter}: human-readable text + \texttt{key=value} extras; JSON mode via \texttt{LOG\_FORMAT=json} \\
\texttt{core/analyzer.py}           & Orchestrates the full pipeline: validate $\to$ VLM $\to$ nutrition $\to$ totals $\to$ persist \\
\texttt{services/ai\_service.py}    & Wraps \texttt{ai.identify\_ingredients} with retry policy and \texttt{asyncio} timeout \\
\texttt{services/nutrition\_cache.py} & In-memory, coroutine-safe TTL cache with \texttt{asyncio.Lock} \\
\texttt{services/retry.py}          & Reusable \texttt{tenacity} \texttt{AsyncRetrying} policy (exponential backoff, jitter) \\
\texttt{concurrency/pipeline.py}    & \texttt{NutritionPipeline}: fans out USDA calls via \texttt{asyncio.gather} + \texttt{Semaphore} \\
\texttt{storage/db.py}              & \texttt{asyncpg} pool creation, schema initialisation, lifespan management \\
\texttt{storage/repository.py}      & SQL CRUD: \texttt{save(AnalysisRecord)}, \texttt{list\_recent(limit, offset)} \\
\texttt{api.py}                     & FastAPI app: \texttt{POST /analyze}, \texttt{GET /history}, \texttt{GET /health}, \texttt{GET /} (web UI) \\
\texttt{cli.py}                     & \texttt{python -m foodanalyzer analyze <path> [{-}{-}offline] [{-}{-}json]} \\
\texttt{static/index.html}          & Single-page browser UI with drag-and-drop upload and analysis history \\
\bottomrule
\end{tabularx}
\end{center}

\textbf{Provider swap test.}
If we switched from Gemini to Anthropic, only \texttt{.env} would change (\texttt{LLM\_PROVIDER} and the API key); no \texttt{src/} files would need editing --- the factory pattern in \texttt{ai/providers/factory.py} handles adapter selection transparently.

\subsection{OOP design \& key abstractions}

The provided \texttt{ai/} package defines an abstract \texttt{VLMProvider} with concrete subclasses (\texttt{GeminiVLM}, \texttt{AnthropicVLM}, \texttt{OpenAIVLM}).
In our SE layer we use \textbf{composition} in \texttt{FoodAnalyzer}: the analyzer receives \texttt{AIService} and \texttt{NutritionPipeline} as constructor arguments, rather than inheriting from any base class.
This is the right tool because the analyzer does not specialise either service --- it coordinates them.
Dependency injection also makes unit tests straightforward: we pass \texttt{FakeVLMProvider} and \texttt{FakeNutritionProvider} without any monkey-patching.

\begin{lstlisting}[style=python, caption={core/analyzer.py (simplified)}]
class FoodAnalyzer:
    """Orchestrates: validate -> VLM -> nutrition -> totals -> persist."""

    def __init__(
        self,
        ai_service: AIService,
        pipeline: NutritionPipeline,
        repo: AnalysisRepository | None = None,
    ) -> None:
        self._ai = ai_service
        self._pipeline = pipeline
        self._repo = repo

    async def analyze(self, image_path: Path) -> AnalysisResult:
        validate_image(image_path)
        ingredients = await self._ai.identify_ingredients(image_path)
        nutrition_map = await self._pipeline.fetch_all(ingredients)
        totals = compute_totals(ingredients, nutrition_map)
        result = AnalysisResult(ingredients=..., totals=totals)
        if self._repo:
            await self._repo.save(AnalysisRecord(...))
        return result
\end{lstlisting}

\textbf{Justification.}
Composition is preferred over inheritance here because \texttt{FoodAnalyzer} is not a specialisation of \texttt{AIService} or \texttt{NutritionPipeline} --- it \emph{uses} them.
A plain class with typed constructor arguments and dependency injection gives us testability, flexibility, and clarity with less machinery than an abstract base class would.

\subsection{Data flow across module boundaries}

Our SE layer defines four Pydantic models that form the typed contract between modules:

\begin{center}\small
\begin{tabularx}{\textwidth}{@{} l X @{}}
\toprule
\textbf{Model} & \textbf{Contents} \\
\midrule
\texttt{IngredientOut}  & Extends \texttt{Ingredient} (from \texttt{ai/schemas.py}) with \texttt{kcal}, \texttt{protein\_g}, \texttt{carbs\_g}, \texttt{fat\_g} \\
\texttt{TotalsOut}      & Aggregated macronutrients across all ingredients \\
\texttt{AnalysisResult} & \texttt{meal\_recognized}, \texttt{image\_path}, \texttt{ingredients: list[IngredientOut]}, \texttt{totals: TotalsOut} \\
\texttt{AnalysisRecord} & Database record: \texttt{id} (UUID), \texttt{image\_path}, \texttt{result} (JSONB), \texttt{created\_at} \\
\bottomrule
\end{tabularx}
\end{center}

No naked dictionaries cross module boundaries.

% =====================================================================
\section{Concurrency \& Performance}\label{sec:concurrency}

\subsection{Why this concurrency model}

We used \textbf{\texttt{asyncio} with \texttt{asyncio.gather}} for the nutrition-lookup fan-out in \texttt{NutritionPipeline}, bounded by \texttt{asyncio.Semaphore(MAX\_PARALLEL)}.
The bottleneck is I/O: each USDA FoodData Central HTTP call takes $\sim$0.5\,s of network round-trip time.
\texttt{asyncio} (single-threaded cooperative concurrency) is the right tool here because the work is I/O-bound, not CPU-bound; threading or multiprocessing would add overhead without benefit.
A \texttt{ProcessPoolExecutor} would be appropriate only if parsing large JSON payloads were a bottleneck, which it is not.

The semaphore bound of 10 was chosen to stay within the USDA API's rate limit (1,000 requests per hour per key) while allowing a realistic meal (5--10 ingredients) to complete in a single round of concurrent calls.
Setting the semaphore to 1 degrades to sequential behaviour; setting it to 100 risks hitting the provider rate limit for large batch requests.

\subsection{Sequential-vs-concurrent benchmark}

Benchmark conditions: simulated USDA latency of 0.5\,s per call (representative of real network RTT); caches cleared before each run; Python~3.13.7; Windows~11.
Command: \texttt{python artefacts/bench\_script.py}.

\begin{center}\small
\begin{tabular}{lcccc}
\toprule
\textbf{Workload} & $N$ & \textbf{Sequential (s)} & \textbf{Concurrent (s)} & \textbf{Speedup} \\
\midrule
Nutrition lookups & 3 & 1.51 & 0.50 & 3.0$\times$ \\
Nutrition lookups & 5 & 2.50 & 0.53 & \textbf{4.8$\times$} \\
Nutrition lookups & 7 & 3.51 & 0.52 & 6.7$\times$ \\
\bottomrule
\end{tabular}
\end{center}

The speedup approaches $N\times$ and slightly exceeds it at $N=7$ due to Python's async scheduling overhead averaging out.
The new bottleneck after parallelisation is the Gemini VLM call (sequential by design --- one image, one response) and the single DB write at the end.

\subsection{Rate limit \& politeness}

We implement \textbf{semaphore-bounded concurrency} as the primary rate-limit strategy: \texttt{Semaphore(MAX\_PARALLEL)} (default 10) ensures we never have more than 10 outstanding USDA requests at once.
In addition, the \texttt{tenacity} retry policy sleeps \texttt{wait\_exponential(multiplier=1, min=2, max=30)} between retries, which acts as a natural back-off on 429 responses.
The configured burst ceiling is 10 requests; a typical 5-ingredient meal analysis makes exactly 5 concurrent USDA calls.

% =====================================================================
\section{Robustness \& Error Handling}\label{sec:robustness}

\subsection{Wrapping the AI module}

Every call to \texttt{ai.*} is wrapped in \texttt{AIService} or \texttt{NutritionPipeline}:

\begin{itemize}[topsep=4pt,itemsep=3pt]
  \item \textbf{Retries.}
    \texttt{tenacity.AsyncRetrying} with \texttt{retry\_if\_exception\_type}, \texttt{wait\_exponential(min=2, max=30)}, \texttt{stop\_after\_attempt(4)}.
    The policy is defined once in \texttt{services/retry.py} and reused across all services.

  \item \textbf{Timeout.}
    \texttt{asyncio.wait\_for(coroutine, timeout=60.0)} in \texttt{AIService.identify\_ingredients};
    20\,s per ingredient lookup in \texttt{NutritionPipeline}.

  \item \textbf{Structured logging.}
    Before each VLM call: \texttt{INFO} with \texttt{image\_path} and \texttt{provider}.
    After success: \texttt{INFO} with \texttt{ingredient\_count} and \texttt{elapsed\_ms}.
    On retry: \texttt{WARNING} with \texttt{attempt} and \texttt{exc\_type}.
    Full payloads available at \texttt{DEBUG} level.

  \item \textbf{Input validation.}
    \texttt{validate\_image(path)} checks file existence, magic bytes
    (JPEG: \texttt{FF D8 FF}, PNG: \texttt{89 50 4E 47}), and size limit before any API call.

  \item \textbf{Response validation.}
    We additionally check \texttt{estimated\_grams > 0} and that the ingredient list is non-empty
    before passing downstream.
\end{itemize}

\subsection{Failure-mode analysis --- one concrete incident}

\begin{center}\small
\begin{tabularx}{\textwidth}{@{} l X @{}}
\toprule
\textbf{Step} & \textbf{Detail} \\
\midrule
Failure injected  & Wrong Gemini image upload API: \texttt{Files.upload()} not supported in \texttt{google-genai==0.3.0} \\
Root cause        & \texttt{ai/providers/google.py} called \texttt{self.\_client.files.upload(file=str(path))}, a v1 API unavailable in 0.3.0 \\
System behaviour  & \texttt{tenacity} retried 4 times, then re-raised as \texttt{AnalysisError}; FastAPI returned HTTP~500 \\
User saw          & \texttt{\{"detail": "Gemini call failed: \ldots unexpected keyword argument 'file'"\}} --- no traceback \\
Fix applied       & Switched to \texttt{types.Part.from\_bytes(data=image\_bytes, mime\_type=mime\_type)} \\
\bottomrule
\end{tabularx}
\end{center}

Log output during the incident:
\begin{lstlisting}[style=bash]
WARNING  attempt=1 exc_type=TypeError retrying_in=2.0s
WARNING  attempt=2 exc_type=TypeError retrying_in=4.0s
WARNING  attempt=3 exc_type=TypeError retrying_in=8.0s
ERROR    ai_service failed after 4 attempts exc=TypeError(...)
\end{lstlisting}

\subsection{Input validation summary}

\begin{center}\small
\begin{tabularx}{\textwidth}{@{} l X l @{}}
\toprule
\textbf{Entry point} & \textbf{Rule} & \textbf{Error on failure} \\
\midrule
\texttt{POST /analyze} & JPEG/PNG magic bytes only & HTTP 422 ``Unsupported image format'' \\
\texttt{POST /analyze} & File $\le$ \texttt{MAX\_IMAGE\_SIZE\_MB} (5\,MB) & HTTP 422 ``Image too large'' \\
\texttt{GET /history}  & \texttt{1 $\le$ limit $\le$ 100} via \texttt{Query(ge=1, le=100)} & HTTP 422 (FastAPI auto) \\
CLI \texttt{analyze}   & File must exist; same size/format checks & \texttt{FileNotFoundError} + message \\
Config / env           & \texttt{pydantic-settings} type validation at startup & \texttt{ValidationError} to stderr \\
\bottomrule
\end{tabularx}
\end{center}

% =====================================================================
\section{Storage \& Caching}\label{sec:storage}

\subsection{Database schema}

We use PostgreSQL via \texttt{asyncpg}.
The schema is initialised by \texttt{storage/db.py:init\_schema()} on application startup:

\begin{lstlisting}[style=sql]
CREATE TABLE IF NOT EXISTS analyses (
    id          UUID        PRIMARY KEY DEFAULT gen_random_uuid(),
    image_path  TEXT        NOT NULL,
    result      JSONB       NOT NULL,
    created_at  TIMESTAMPTZ NOT NULL DEFAULT now()
);
CREATE INDEX IF NOT EXISTS analyses_created_at_idx
    ON analyses (created_at DESC);
\end{lstlisting}

The full \texttt{AnalysisResult} (ingredients + totals) is stored as JSONB in the \texttt{result} column.
This avoids the need for a separate \texttt{ingredients} table and makes the history query a single \texttt{SELECT} with no joins.
The \texttt{image\_path} field stores the relative path under \texttt{uploads/}; raw image bytes live on the local filesystem.

\subsection{Caching policy}

\texttt{NutritionCache} in \texttt{services/nutrition\_cache.py} is an in-memory, per-process TTL cache for USDA nutrition lookups.

\begin{center}\small
\begin{tabular}{ll}
\toprule
\textbf{Property} & \textbf{Value} \\
\midrule
Key        & Case-insensitive ingredient name \\
TTL        & \texttt{NUTRITION\_CACHE\_TTL\_SECONDS} (default 86400\,s = 24\,h) \\
Eviction   & Lazy (checked on read) \\
Safety     & \texttt{asyncio.Lock} prevents concurrent writer races \\
Hit rate   & $> 60\%$ on a run over 16 sample images (shared ingredients: rice, chicken, broccoli) \\
\bottomrule
\end{tabular}
\end{center}

Staleness risk: USDA periodically updates nutrient values.
A 24-hour TTL is a reasonable trade-off for a demo system; a production deployment would use a persistent Redis cache keyed on the USDA food ID.

% =====================================================================
\section{Testing}\label{sec:testing}

\subsection{Strategy \& coverage}

All 158 tests run \textbf{offline} --- no network, no database, no real API keys required.
The AI module is mocked with \texttt{FakeVLMProvider} and \texttt{FakeNutritionProvider} defined in \texttt{tests/conftest.py}.
HTTP calls to USDA are intercepted by \texttt{unittest.mock.AsyncMock}.
The PostgreSQL pool is replaced by an in-memory mock repository.

\begin{center}\small
\begin{tabular}{lc}
\toprule
\textbf{Suite} & \textbf{Result} \\
\midrule
\texttt{src/foodanalyzer/} (SE layer) & \textbf{79\%} coverage \\
Provided \texttt{tests/test\_ai\_smoke.py} & \textbf{29/29 passing} \\
Total tests collected & \textbf{158 passed} \\
\bottomrule
\end{tabular}
\end{center}

The main uncovered lines are in \texttt{cli.py} (CLI integration paths) and a few lifespan edge-cases in \texttt{api.py}.

\subsection{Notable tests}

\begin{enumerate}[topsep=4pt,itemsep=4pt]
  \item \textbf{Happy-path end-to-end} --- \texttt{test\_api.py::test\_analyze\_success}\\
    Posts a real PNG to \texttt{POST /analyze} via FastAPI's \texttt{TestClient}, with mocked services.
    Asserts HTTP~200, \texttt{meal\_recognized=True}, non-empty ingredient list, and \texttt{totals.kcal > 0}.
    Catches any regression in request parsing, response serialisation, or pipeline wiring.

  \item \textbf{Concurrency failure degradation} --- \texttt{test\_concurrency.py::test\_one\_task\_raises\_others\_complete}\\
    Creates a \texttt{NutritionPipeline} with five ingredients, one of which raises \texttt{RuntimeError}.
    Asserts that the pipeline returns partial results for the remaining four and logs a warning.
    This test exposed a design issue: the original implementation cancelled all tasks on the first failure;
    we fixed it to use \texttt{asyncio.gather(return\_exceptions=True)} and filter.

  \item \textbf{Magic-byte validation} --- \texttt{test\_validation.py::test\_wrong\_magic\_bytes}\\
    Submits a \texttt{.jpg} file with PDF magic bytes (\texttt{\%PDF}).
    Asserts that \texttt{validate\_image\_bytes} raises \texttt{ValueError("Unsupported image format")} before any API call.
    This caught a real bug: the original code checked only the file extension.
\end{enumerate}

% =====================================================================
\section{Deployment \& Reproducibility}\label{sec:deploy}

\subsection{Docker image}

\begin{lstlisting}[style=bash]
# Build
docker build -t food-analyzer .

# Run (API server on port 8000)
docker run -p 8000:8000 --env-file .env food-analyzer

# Full dev stack: PostgreSQL + pgAdmin + app
docker compose up -d
\end{lstlisting}

The image is built on \textbf{\texttt{python:3.12-slim}} (Debian Bookworm slim).
A non-root user (\texttt{appuser}) is created and used at runtime.
The application starts with \texttt{CMD ["uvicorn", "foodanalyzer.api:app", "--host", "0.0.0.0", "--port", "8000"]}.
Port 8000 is declared with \texttt{EXPOSE 8000}.
Image size is approximately 350--400\,MB including all Python dependencies; a fresh \texttt{python:3.12-slim} is $\sim$120\,MB, so our dependencies add $\sim$250\,MB.

\subsection{Environment variables}

\begin{center}\small
\begin{tabularx}{\textwidth}{@{} l l X @{}}
\toprule
\textbf{Variable} & \textbf{Default} & \textbf{Effect if missing} \\
\midrule
\texttt{LLM\_PROVIDER}      & \texttt{anthropic}  & Falls back to Anthropic; set to \texttt{gemini} for Google \\
\texttt{GOOGLE\_API\_KEY}   & ---                 & VLM calls fail with auth error (retried, then 500 returned) \\
\texttt{USDA\_API\_KEY}     & ---                 & Nutrition lookups fail; totals show 0 (graceful degradation) \\
\texttt{POSTGRES\_HOST}     & \texttt{localhost}  & App starts; first DB write logs error and returns 500 \\
\texttt{POSTGRES\_PASSWORD} & \texttt{dev}        & Must match Docker Compose value; mismatch fails connection \\
\texttt{MAX\_PARALLEL}      & \texttt{10}         & Controls USDA concurrency; safe to leave at default \\
\texttt{LOG\_LEVEL}         & \texttt{INFO}       & Set to \texttt{DEBUG} for full payload logging \\
\bottomrule
\end{tabularx}
\end{center}

% =====================================================================
\section{Limitations \& Future Work}\label{sec:limits}

\begin{itemize}[topsep=4pt,itemsep=4pt]
  \item \textbf{In-memory cache is per-process and non-persistent.}
    Restarting the server clears the nutrition cache.
    A Redis-backed cache would survive restarts and be shared across multiple replicas.

  \item \textbf{CLI test coverage is 0\,\%.}
    We manually verified CLI functionality but did not write subprocess-level tests.
    A \texttt{typer.testing.CliRunner} harness would close this gap.

  \item \textbf{Image storage is local filesystem.}
    In a multi-replica deployment, uploaded images would not be shared across instances.
    Moving to an object store (S3, GCS) is the obvious next step.

  \item \textbf{No authentication on the API.}
    \texttt{POST /analyze} is open.
    In production, per-client rate limiting and API-key authentication would be required.

  \item \textbf{USDA search is text-matching dependent.}
    If Gemini returns an ingredient name that USDA does not recognise (e.g.\ a regional dish name),
    the nutrition lookup silently returns zero values.
    A fuzzy-match or category-fallback strategy would improve accuracy.
\end{itemize}

% =====================================================================
\section{Tools \& Acknowledgements}\label{sec:tools}

Throughout the project we used \textbf{Claude}  as a coding assistant in the same way a developer would use it in industry: for drafting boilerplate, suggesting library usage, and helping diagnose bugs.
All AI-generated suggestions were reviewed, tested, and adapted by the team before being committed.
In particular, the assistant helped us explore \texttt{tenacity} retry configuration, \texttt{asyncio.gather} error-handling patterns, and diagnose the \texttt{pydantic-settings} vs \texttt{os.getenv} environment-loading incompatibility.

Every team member understands the full codebase and is able to explain and defend every line of code during the oral defense.
``The AI wrote it'' is not an answer any of us will use.

% =====================================================================
\section*{References}
\addcontentsline{toc}{section}{References}

\begin{enumerate}[label={[\arabic*]},leftmargin=2.4em,itemsep=3pt]
  \item Google Gemini API documentation. \url{https://ai.google.dev/docs}
  \item USDA FoodData Central API. \url{https://fdc.nal.usda.gov/api-guide}
  \item \texttt{tenacity} retry library. \url{https://tenacity.readthedocs.io/}
  \item FastAPI documentation. \url{https://fastapi.tiangolo.com/}
  \item \texttt{pydantic-settings} documentation. \url{https://docs.pydantic.dev/latest/concepts/pydantic_settings/}
  \item \texttt{asyncpg} documentation. \url{https://magicstack.github.io/asyncpg/}
  \item Docker --- Best practices for writing Dockerfiles. \url{https://docs.docker.com/develop/develop-images/dockerfile_best-practices/}
\end{enumerate}

\appendix
% =====================================================================
\section{Reproducing the benchmark}

Anyone with the repository and a valid \texttt{.env} file can reproduce the concurrency benchmark from Section~\ref{sec:concurrency}:

\begin{lstlisting}[style=bash]
git clone https://github.com/AMammedova/food-analyzer-byte-bite.git
cd food-analyzer-byte-bite
cp .env.example .env          # fill in GOOGLE_API_KEY and USDA_API_KEY

python -m venv .venv
.venv\Scripts\activate        # Windows
pip install -r requirements.txt

$env:PYTHONPATH="src"         # Windows PowerShell
python artefacts/bench_script.py
\end{lstlisting}

Expected output (simulated 0.5\,s USDA latency per call):

\begin{lstlisting}[style=bash]
Benchmark: simulated USDA latency 0.5s per call
------------------------------------------------------------
N= 3  Sequential=1.51s  Concurrent=0.50s  Speedup=3.0x
N= 5  Sequential=2.50s  Concurrent=0.53s  Speedup=4.8x
N= 7  Sequential=3.51s  Concurrent=0.52s  Speedup=6.7x
------------------------------------------------------------
\end{lstlisting}

\vspace{1em}\hrule\vspace{0.6em}
\noindent\small\textit{End of report --- Byte Bite team, AI Academy Spring 2026.}

\end{document}
